\section{ภาพรวมระบบ}
ระบบคำร้องขอลานิสิตคณะวิศวกรรมศาสตร์ใช้สำหรับสร้างคำร้อง ติดตามผล พิจารณาอนุมัติ รับทราบผลกระทบต่อรายวิชา และจัดทำรายงาน โดยเชื่อมข้อมูลบัญชี บุคลากร นิสิต รายวิชา และโครงสร้างภาควิชาจากบริการที่เกี่ยวข้อง

\subsection{กลุ่มผู้ใช้งาน}

\begin{table}[H]
  \centering
  \begingroup
  \renewcommand{\arraystretch}{1.22}
  \setlength{\tabcolsep}{7pt}
  \begin{tabularx}{\textwidth}{|>{\ttfamily\fontsize{13pt}{16pt}\selectfont\raggedright\arraybackslash}p{4.6cm}|>{\raggedright\arraybackslash}p{3.1cm}|>{\raggedright\arraybackslash}X|}
    \hline
    \multicolumn{1}{|l|}{\tableheader{Role code}}
    & \tableheader{ชื่อบทบาท}
    & \tableheader{หน้าที่หลัก} \\\hline

    STUDENT
    & นิสิต
    & สร้างคำร้อง ดูสถานะ เปิดเอกสาร และตรวจรายวิชาของตน \\\hline
ADVISOR
    & อาจารย์ที่ปรึกษา
    & พิจารณาคำร้องในขั้นอาจารย์ที่ปรึกษา \\\hline
DEPARTMENT\_HEAD
    & หัวหน้าภาควิชา
    & พิจารณาคำร้องในขั้นหัวหน้าภาควิชา \\\hline
ASSOCIATE\_DEAN
    & รองคณบดี
    & พิจารณาคำร้องที่ลาเกิน 15 วัน \\\hline
LECTURER
    & อาจารย์ประจำวิชา
    & เปิดอ่านและรับทราบคำร้องที่อนุมัติแล้วในรายวิชาที่สอน \\\hline
ADMIN\_STAFF
    & เจ้าหน้าที่
    & ตรวจรายการ จัดการภาคการศึกษา โครงสร้าง และตั้งค่าระบบ \\\hline
  \end{tabularx}
  \endgroup
  \caption{บทบาทและหน้าที่หลักในระบบ}
\end{table}

หนึ่งบัญชีอาจมีหลายบทบาท ระบบจะแสดงเมนูตามสิทธิ์ที่ได้รับ

\subsection{กระบวนการโดยสรุป}
\begin{enumerate}
  \item นิสิตกรอกคำร้องและเลือกรายวิชาที่ได้รับผลกระทบ
  \item อาจารย์ที่ปรึกษาพิจารณา
  \item หัวหน้าภาควิชาพิจารณา
  \item รองคณบดีพิจารณาเพิ่มเติมเฉพาะกรณีลามากกว่า 15 วัน
  \item เมื่ออนุมัติครบ อาจารย์ประจำวิชาที่เกี่ยวข้องเปิดอ่านและรับทราบ
\end{enumerate}
หากผู้พิจารณาขั้นใดไม่อนุมัติ คำร้องจะสิ้นสุดและไม่ส่งต่อไปยังขั้นถัดไป

\subsection{ขอบเขตเอกสาร}
บทที่ 2--4 เป็นคู่มือสำหรับผู้พัฒนา บทที่ 5 อธิบายกฎของ workflow และบทที่ 6--9 เป็นคู่มือใช้งานแยกตามบทบาท รวมถึงการแก้ปัญหาและความปลอดภัย
